\documentclass[11pt]{article}
\usepackage{mathblatt}
\begin{document}
\weit
\typeout{PROBE-BEGIN ableitung-und-aenderungsrate-e1-k1-s2-v2}
\begin{aufgabe}{\detokenize{ableitung-und-aenderungsrate-e1-k1-s2-v2}}
\begin{teile}
\teil Probe

\begin{ksys}[xmin=0,xmax=4,ymin=0,ymax=1800,xstep=1,ystep=200,xlabel=t in h,ylabel=a,ablesen]\funktion{100*\x^2}{a}\end{ksys}
\end{teile}
\end{aufgabe}
\typeout{PROBE-END ableitung-und-aenderungsrate-e1-k1-s2-v2}
\typeout{PROBE-BEGIN ableitungsregeln-e2-k1-s5-v1}
\begin{aufgabe}{\detokenize{ableitungsregeln-e2-k1-s5-v1}}
\begin{teile}
\teil Probe

\begin{ksys}[xmin=-3,xmax=4,ymin=-3,ymax=5,ablesen] \funktion{0.5*\x^2-1}{f} \parabel{1}{1}{-2}{g} \tiefpunkt{1}{-2}{T} \end{ksys}
\end{teile}
\end{aufgabe}
\typeout{PROBE-END ableitungsregeln-e2-k1-s5-v1}
\typeout{PROBE-BEGIN bedingte-wahrscheinlichkeit-und-bayes-e3-k1-s3-v1}
\begin{aufgabe}{\detokenize{bedingte-wahrscheinlichkeit-und-bayes-e3-k1-s3-v1}}
\begin{teile}
\teil Probe

\begin{ksys}[xmin=0,xmax=1,ymin=0,ymax=1,xstep=0.2,ystep=0.2,ablesen] \funktion{0.15/(0.15+0.7*\x)}{A} \funktion{1-\x}{B} \funktion{0.15/(0.15+0.7*(1-\x))}{C} \end{ksys}
\end{teile}
\end{aufgabe}
\typeout{PROBE-END bedingte-wahrscheinlichkeit-und-bayes-e3-k1-s3-v1}
\typeout{PROBE-BEGIN extremalprobleme-e1-k1-s2-v1}
\begin{aufgabe}{\detokenize{extremalprobleme-e1-k1-s2-v1}}
\begin{teile}
\teil Probe

\begin{ksys}[xmin=-5,xmax=5,ymin=-5,ymax=5] \funktion{4-0.25*\x^2}{f} \funktion{-4+0.25*\x^2}{} \end{ksys}
\end{teile}
\end{aufgabe}
\typeout{PROBE-END extremalprobleme-e1-k1-s2-v1}
\typeout{PROBE-BEGIN flaecheninhalt-durch-integration-e1-k1-s0-v1}
\begin{aufgabe}{\detokenize{flaecheninhalt-durch-integration-e1-k1-s0-v1}}
\begin{teile}
\teil Probe

\begin{ksys}[xmin=-3,xmax=4,ymin=-5,ymax=5,ablesen]\funktion{0.5*\x^2+1}{f}\end{ksys}
\end{teile}
\end{aufgabe}
\typeout{PROBE-END flaecheninhalt-durch-integration-e1-k1-s0-v1}
\end{document}
